
%%%%%%%%%%%%%%%%%%%%%%%%%%%%%%%%%%%%%%%%
%           main body
%%%%%%%%%%%%%%%%%%%%%%%%%%%%%%%%%%%%%%%%


%-----------------------------------
%	TITLE PAGE
%-----------------------------------

\title[Convex Optimization]{\LARGE \bfseries 第四部分\quad 凸优化} % The short title appears at the bottom of every slide, the full title is only on the title page
\author[Newton’s Method and Interior-Point Methods] % Your institution as it will appear on the bottom of every slide, may be shorthand to save space
{\Large \bfseries \S 5. 牛顿法与内点法}
% \author{Singular Value Decomposition} % Your name
% \date{\today} % Date, can be changed to a custom date
\date{}

%------------------------------------------------

\begin{document}


%------------------------------------------------

\begin{frame}

\titlepage % Print the title page as the first slide

\end{frame}

%------------------------------------------------

\begin{frame}
\frametitle{内容提要} % Table of contents slide, comment this block out to remove it
\tableofcontents % Throughout your presentation, if you choose to use \section{} and \subsection{} commands, these will automatically be printed on this slide as an overview of your presentation
\end{frame}

%-----------------------------------
%	PRESENTATION SLIDES
%-----------------------------------

\section{无约束优化}

%------------------------------------------------

\begin{frame}

\begin{block}{无约束优化}
无约束凸优化问题具有以下形式
$$\text{minimize} \quad f(\mathbf{x})$$
由于$f$是凸函数，那么
$$\mathbf{x}^* \text{ 是最优解 } \Longleftrightarrow \nabla f(\mathbf{x}^*) = \mathbf{0}$$
这个方程（组）一般没有解析解，必须采用迭代算法求解，即寻找所谓的极小化点列
$$\mathbf{x}^{(0)}, \ldots, \mathbf{x}^{(k)}, \ldots,$$
使得当$k\to\infty$时，$f(\mathbf{x}^{(k)}) \to p^*$
\end{block}

\end{frame}

%------------------------------------------------

\begin{frame}

\begin{block}{下降方法}
极小化点列将按以下方式产生
$$\mathbf{x}^{(k+1)} = \mathbf{x}^{(k)} + t^{(k)} \Delta \mathbf{x}^{(k)}, \text{ 且有 } f(\mathbf{x}^{(k+1)}) < f(\mathbf{x}^{(k)})$$
其中
\begin{itemize}
\item $\Delta \mathbf{x}^{(k)}$被称为第$k$次迭代搜索方向，或步径
\item $t^{(k)} > 0$被称作第$k$次迭代步长
\end{itemize}
\end{block}

\end{frame}

%------------------------------------------------

\begin{frame}

\begin{algorithm}[H]
\SetAlgoLined
\DontPrintSemicolon
\SetKwInOut{Input}{Input}
\SetKwInOut{Output}{Output}
\SetKwRepeat{Do}{do}{while}

给定初始点$\mathbf{x}$ \;
\Do{达到停止条件}
{
确定下降方向$\Delta \mathbf{x}$ \;
搜索并选择步长$t$ \;
$ \mathbf{x} \gets \mathbf{x} + t \Delta \mathbf{x}$ \;
}
\Output{近似最优解$\mathbf{x}$}
\caption{通用下降方法}
\end{algorithm}

\vspace{1em}
\pause

\begin{itemize}
\item $\Delta \mathbf{x} = - \nabla f(\mathbf{x}) \; \rightsquigarrow \; \text{\color{red} 梯度下降法}$
\item $\Delta \mathbf{x}_{nsd} = \argmin\limits_{\mathbf{v}} \{ \nabla f(\mathbf{x})^T \mathbf{v} \ |\ \Vert \mathbf{v} \Vert = 1 \} \; \rightsquigarrow \; \text{\color{red} 最速下降法}$
\item $\Delta \mathbf{x}_{nt} = - H(f)(\mathbf{x})^{-1} \nabla f(\mathbf{x}) \; \rightsquigarrow \; \text{\color{red} Newton法}$
\end{itemize}

\end{frame}

%------------------------------------------------

\begin{frame}

\begin{block}{Newton法}
\begin{itemize}
\item 函数$f$在$\mathbf{x}$处的二阶Taylor近似为
$$f(\mathbf{x} + \mathbf{v}) = f(\mathbf{x}) + \nabla f(\mathbf{x})^T \mathbf{v} + \dfrac{1}{2} \mathbf{v}^T H(f)(\mathbf{x}) \mathbf{v}$$
$\mathbf{x} + \Delta \mathbf{x}_{nt}$ 极小化了$f$在$\mathbf{x}$处的二阶近似
\vspace{1em}
\item Newton法的步径（或搜索方向）也是$\mathbf{x}$处采用正定阵$H(f)(\mathbf{x})$定义的（内积导出的）范数
$$\Vert \mathbf{u} \Vert_{H(f)(\mathbf{x})} := (\mathbf{u}^T H(f)(\mathbf{x}) \mathbf{u})^{1/2}$$
对应的最速下降法的步径。
\end{itemize}
\end{block}

\end{frame}

%------------------------------------------------

\begin{frame}

\begin{block}{Newton法}
\begin{itemize}
\item 定义$\mathbf{x}$处的Newton减量（Decrement）
\begin{align*}
\lambda(\mathbf{x}) & = (\Delta \mathbf{x}_{nt}^T H(f)(\mathbf{x}) \Delta \mathbf{x}_{nt})^{1/2} \\
& = (\nabla f(\mathbf{x})^T H(f)(\mathbf{x})^{-1} \nabla f(\mathbf{x}))^{1/2}
\end{align*}
可以用于设计迭代停止准则
\end{itemize}
\end{block}

\end{frame}

%------------------------------------------------

\begin{frame}

\begin{algorithm}[H]
\SetAlgoLined
\DontPrintSemicolon
\SetKwInOut{Input}{Input}
\SetKwInOut{Output}{Output}
\SetKwRepeat{Do}{do}{while}

给定初始点$\mathbf{x}$, 以及误差阈值$\varepsilon$ \;
\Do{1}
{
计算Newton步径$\Delta \mathbf{x}_{nt} = - H(f)(\mathbf{x})^{-1} \nabla f(\mathbf{x})$, 以及Newton减量平方$\lambda(\mathbf{x})^2 = \nabla f(\mathbf{x})^T H(f)(\mathbf{x})^{-1} \nabla f(\mathbf{x})$\;
如果$\lambda(\mathbf{x})^2 \le \varepsilon$, 退出循环 \;
搜索并选择步长$t$ \;
$ \mathbf{x} \gets \mathbf{x} + t \Delta \mathbf{x}_{nt}$ \;
}
\Output{近似最优解$\mathbf{x}$}
\caption{Newton法}
\end{algorithm}

\end{frame}

%------------------------------------------------

\section{等式约束的Newton法}

%------------------------------------------------

\begin{frame}

\begin{block}{等式约束的Newton法}
考虑如下等式约束的凸优化问题
\vspace{-0.5em}
\begin{eqnarray*}
& \text{minimize} & f(\mathbf{x}) \\
& \text{subject to} & A\mathbf{x} = \mathbf{b}
\end{eqnarray*}
其中目标函数$f$是二次连续可微的凸函数。

\pause
\vspace{1em}

为得到其在可行点$\mathbf{x}$处的Newton步径$\Delta \mathbf{x}_{nt}$, 将目标函数$f$替换为$f$在$\mathbf{x}$附近的二阶Taylor近似，从而得到如下的关于变量$\mathbf{v}$的等式约束的二次凸问题
\vspace{-0.5em}
\begin{eqnarray*}
& \text{minimize} & g(\mathbf{v}) = f(\mathbf{x}) + \nabla f(\mathbf{x})^T \mathbf{v} + \dfrac{1}{2} \mathbf{v}^T H(f)(\mathbf{x}) \mathbf{v} \\
& \text{subject to} & A(\mathbf{x} + \mathbf{v}) = \mathbf{b} \quad (\Longrightarrow \; A\mathbf{v} = \mathbf{0})
\end{eqnarray*}
关键就在于Newton方向不但要是下降的方向，还要是可行的（满足约束条件）
\end{block}

\end{frame}

%------------------------------------------------

\begin{frame}

\begin{block}{等式约束的Newton法}
以上的等式约束的二次凸问题的最优解$\mathbf{v}^* = \Delta \mathbf{x}_{nt}$由KKT矩阵决定的线性方程组决定
\[
\begin{pmatrix} H(f)(\mathbf{x}) & A^T \\ A & 0 \end{pmatrix} \begin{pmatrix} \Delta \mathbf{x}_{nt} \\ \boldsymbol\omega \end{pmatrix} = \begin{pmatrix} -\nabla f(\mathbf{x}) \\ \mathbf{0} \end{pmatrix}
\]
剩下的工作就和无约束情况下Newton法类似了。
\end{block}

\end{frame}

%------------------------------------------------

\begin{frame}

\begin{algorithm}[H]
\SetAlgoLined
\DontPrintSemicolon
\SetKwInOut{Input}{Input}
\SetKwInOut{Output}{Output}
\SetKwRepeat{Do}{do}{while}

给定可行初始点$\mathbf{x}$, 以及误差阈值$\varepsilon$ \;
\Do{1}
{
通过解KKT方程组计算Newton步径$\Delta \mathbf{x}_{nt}$, 以及相应的Newton减量平方$\lambda(\mathbf{x})^2 = \Delta \mathbf{x}_{nt}^T H(f)(\mathbf{x}) \Delta \mathbf{x}_{nt}$\;
如果$\lambda(\mathbf{x})^2 \le \varepsilon$, 退出循环 \;
搜索并选择步长$t$ \;
$ \mathbf{x} \gets \mathbf{x} + t \Delta \mathbf{x}_{nt}$ \;
}
\Output{近似最优解$\mathbf{x}$}
\caption{等式约束问题的Newton法}
\end{algorithm}

\end{frame}

%------------------------------------------------

\section{内点法}

%------------------------------------------------

\begin{frame}

\begin{block}{不等式约束的凸优化问题}
考虑含有不等式约束的凸优化问题
\begin{eqnarray*}
& \text{minimize} & f_0(\mathbf{x}) \\
& \text{subject to} & f_i(\mathbf{x}) \le 0, \quad i = 1,\ldots,m \\
& & A\mathbf{x} = \mathbf{b}
\end{eqnarray*}
其中
\begin{itemize}
\item $f_0,f_1,\ldots,f_m: \mathbb{R}^n \to \mathbb{R}$为二次连续可微的凸函数
\item $A \in M_{p\times n}(\mathbb{R})$ 且 $r(A) = p$
\item 最优值$p^*$有限且能取到
\item 问题是严格可行的，即（Slater准则）
$$\exists \; \mathbf{x} \in \operatorname{relint} \mathcal{D}, \; \text{ 使得 }\; 
\begin{array}{l}
f_i(\mathbf{x}) < 0, \; i = 1,\ldots,m \\
A\mathbf{x} = \mathbf{b}
\end{array}$$
\end{itemize}
\end{block}

\end{frame}

%------------------------------------------------

\begin{frame}

\begin{block}{解决问题的思路}
将原问题近似转换为（一系列）等式约束问题，然后应用Newton法求解。
\end{block}

\end{frame}

%------------------------------------------------

\begin{frame}

\begin{block}{消去不等式条件}
我们将不等式约束隐含在目标函数中，把原问题重新表述为如下的等式约束问题
\begin{eqnarray*}
& \text{minimize} & f_0(\mathbf{x}) + \sum\limits_{i=1}^m I_-(f_i(\mathbf{x})) \\
& \text{subject to} & A\mathbf{x} = \mathbf{b}
\end{eqnarray*}
其中$I_-: \mathbb{R} \to \mathbb{R}\cup\{\infty\}$是非正实数的示性函数：
\[
I_-(u) = \begin{cases}
0, & u \le 0 \\
\infty, & u > 0 
\end{cases}
\]

\vspace{1em}
\pause

但是以上等式约束问题的目标函数不可微，不能直接用Newton法
\end{block}

\end{frame}

%------------------------------------------------

\begin{frame}

\begin{block}{对数障碍，Logarithmic Barrier}
我们将用对数函数来近似示性函数$I_-(u)$
$$\widehat{I}_-(u) = -\dfrac{1}{t}\ln(-u), \quad u < 0$$
\vspace{-1em}
\begin{itemize}
\item $t>0$是确定近似精度的参数，$t$越大，近似精度越高
\item $\widehat{I}_-(u)$是关于$u$的凸的可微增函数
\item $u\to 0-$时，$\widehat{I}_-(u)\to +\infty$
\end{itemize}

\vspace{1em}

令
$$\phi(\mathbf{x}) := \sum\limits_{i=1}^m \left( -\ln(-f_i(\mathbf{x})) \right)$$
此函数被称作是原问题的对数障碍函数，其定义域是满足原问题中严格不等式约束的点集。
\end{block}

\end{frame}

%------------------------------------------------

\begin{frame}

\begin{block}{对数障碍，Logarithmic Barrier}
于是，我们得到了如下的近似问题
\begin{eqnarray*}
& \text{minimize} & f_0(\mathbf{x}) + \sum\limits_{i=1}^m \left( -\dfrac{1}{t}\ln(-f_i(\mathbf{x})) \right) \\
& \text{subject to} & A\mathbf{x} = \mathbf{b}
\end{eqnarray*}
或等价地（把下述带对数障碍的凸优化问题记作$S(t)$）
\begin{eqnarray*}
& \text{minimize} & tf_0(\mathbf{x}) + \phi(\mathbf{x}) \\
& \text{subject to} & A\mathbf{x} = \mathbf{b}
\end{eqnarray*}
此时的目标函数是可微凸函数，可以用Newton法来解决问题。
\end{block}

\end{frame}

%------------------------------------------------

\begin{frame}

\begin{block}{中心路径，Central Path}
\begin{itemize}
\item 对任意$t>0$, 令$\mathbf{x}^*(t)$为以上近似问题的最优解，被称作一个中心点
\item $\{\mathbf{x}^*(t) \ |\ t > 0\}$被称作原问题的中心路径
\item {\bfseries 重要结论}：每个中心点$\mathbf{x}^*(t)$（在目标函数上的取值）是和原问题的最优值$p^*$之间的偏差不超过$\dfrac{m}{t}$的次优解，即
$$f_0(\mathbf{x}^*(t)) - p^* \le \dfrac{m}{t}$$
也就是说，当$t\to+\infty$, $\mathbf{x}^*(t)$收敛于最优解
\end{itemize}
\end{block}

\end{frame}

%------------------------------------------------

\begin{frame}

\begin{algorithm}[H]
\SetAlgoLined
\DontPrintSemicolon
\SetKwInOut{Input}{Input}
\SetKwInOut{Output}{Output}
\SetKwRepeat{Do}{do}{while}

给定严格可行初始点$\mathbf{x}$, $t > 0$, $\mu > 1$, 以及误差阈值$\varepsilon$ \;
\Do{1}
{
中心点步骤（外部迭代）：\newline
\hspace{1em}以$\mathbf{x}$为初始值，在$A\mathbf{x} = \mathbf{b}$的约束条件下用Newton法（Newton法内的循环迭代被称为内部迭代）来解对数障碍问题$S(t)$，最终确定中心点$\mathbf{x}^*(t)$ \;
$\mathbf{x} \gets \mathbf{x}^*(t)$ \;
如果$m/t \le \varepsilon$, 退出循环 \;
$t \gets \mu t$
}
\Output{近似最优解$\mathbf{x}$}
\caption{内点法（障碍方法）}
\end{algorithm}

\end{frame}

%------------------------------------------------

% \begin{frame}

% \begin{block}{Vandermonde行列式的计算}
% \begin{enumerate}
% % \addtocounter{enumi}{0}
% \item 把第$n-1$行乘以$-a_1$加到第$n$行，再把第$n-2$行乘以$-a_1$加到第$n-1$行。按照这种顺序依次向上，直到用第$2$行乘以$-a_1$加到第$1$行。由行列式性质，$V_n$的值不变。
% \[
% V_n = \det \begin{pmatrix}
% 1 & 1 & \cdots & 1 \\
% 0 & a_2-a_1 & \cdots & a_n-a_1 \\
% \vdots & \vdots & & \vdots \\
% 0 & a_2^{n-2}(a_2-a_1) & \cdots & a_n^{n-2}(a_n-a_1) \\
% \end{pmatrix}
% \]
% \end{enumerate}
% \end{block}

% \end{frame}

%------------------------------------------------
% % closing page
% \begin{frame}
% \HUGE{\centerline{\bfseries {\color{gray} The} {\color{zkblue} End}}}

% \pause

% \vspace{1.2em}

% \Large{\centerline{\bfseries {\color{gray}Thank} \quad {\color{zkblue} You}}}

% \end{frame}

%----------------------------------------------------------------------------------------

\end{document}
